\section*{Chapter \ref{ch:qm:state-space-models-and-the-kalman-filter} --- State-Space Models and the Kalman Filter}

\begin{solution}{exo:qm:state-space-models-and-the-kalman-filter:1}
$q = 0.04$: $p = (0.04 + \sqrt{0.0016 + 0.16})/2 = 0.221$ and $K = 0.221/1.221 = 0.181$. Old observations' weights decay like $0.819^k$, which halves in
$\ln\frac12/\ln0.819 = 3.5$ days.
\end{solution}

\begin{solution}{exo:qm:state-space-models-and-the-kalman-filter:2}
Predicted variance $0.01 + 0.0002 = 0.0102$. \omterm{def:qm:state-space-models-and-the-kalman-filter:kf}{Innovation} $v = 1.2 - 0.8 \times 2 = -0.4$, $F = 4 \times 0.0102 + 0.5 = 0.5408$, gain $K = 0.0102 \times 2/0.5408 =
0.0377$. Updated ratio $0.8 - 0.0377 \times 0.4 = 0.785$, variance $0.0102(1 - 0.0377 \times 2) = 0.0094$.
\end{solution}

\begin{solution}{exo:qm:state-space-models-and-the-kalman-filter:3}
Without an observation there is no update, only the prediction step, which adds the state noise: in the \omterm{def:qm:state-space-models-and-the-kalman-filter:ssm}{local level model} the variance grows by $Q$ that day,
instead of first shrinking by the update. The uncertainty reflects that the state kept moving while nothing was seen.
\end{solution}

\begin{solution}{exo:qm:state-space-models-and-the-kalman-filter:4}
For small $q$, $p = \frac12(q + \sqrt{q^2 + 4q}) \approx \sqrt q$ and $K = p/(1 + p) \approx \sqrt q$. The filtered level is $a_{t+1} = (1 - K)a_t + Ky_t$, an EMA with $\lambda
= 1 - K$; for $q = 0.0004$, $K = 0.0198$ and $\lambda = 0.980$.
\end{solution}

\begin{solution}{exo:qm:state-space-models-and-the-kalman-filter:5}
State $\alpha_t = x_t$, $Z = 1$, $T = \phi$, $Q = \Var(\eta)$, $H = \Var(\varepsilon)$. The steady predicted variance solves $P = \phi^2\bigl(P - P^2/(P + H)\bigr) + Q$.
\end{solution}

\begin{solution}{exo:qm:state-space-models-and-the-kalman-filter:6}
$1/(0.25 + 0.09 + 0.01 + 0.0025 + 0.0025) = 1/0.355 = 2.8$ particles' worth out of five.
\end{solution}

\begin{solution}{exo:qm:state-space-models-and-the-kalman-filter:7}
\texttt{tradeoff()}: $Q \times 0.1$: 0.987 of the window's error, half-response 96 days; $\times 1$: 0.993, 30 days; $\times 10$: 1.006, 9.6 days; $\times 100$: 1.062, 3.0 days.
Slower than the likelihood's choice is marginally better here; faster costs noise.
\end{solution}

\begin{solution}{exo:qm:state-space-models-and-the-kalman-filter:8}
The smoothed ratio for day $t$ uses prices after day $t$, so the backtest hedges with information the desk could not have had. Use the filtered ratio from the
previous close; the smoother belongs to research on what the relationship was, not to a simulation of trading.
\end{solution}

\begin{solution}{pb:qm:state-space-models-and-the-kalman-filter:1}
\textbf{1.} The EIA WTI spot price settled at $-36.98$ dollars; WTI changed by $-55.29$ and $+45.89$ on 20 and 21 April while Brent changed by $-2.39$ and $-8.24$.
Both days are treated as missing: the filter skips the update, and the rolling regression and the evaluation skip the days.
\textbf{2.} 0.41 in 2014, 0.94 in 2021 and 1.01 in 2026.
\textbf{3.} WTI's daily volatility ranged from 0.78 dollars (2017) to 3.62 (2026); with constant $H$ the volatile periods dominate the likelihood and push $Q$ up;
scaling makes the observation noise closer to homoskedastic.
\textbf{4.} 3.324 unhedged and 1.329 with the first 60 days' ratio held fixed (dollars squared per barrel per day, 4\,035 days).
\textbf{5.} 1.276.
\textbf{6.} $H = 0.508$, $Q = 0.000238$, \omterm{def:qm:state-space-models-and-the-kalman-filter:snr}{signal-to-noise ratio} 0.00047.
\textbf{7.} $\ln2/\sqrt{0.00047 \times 1.12} = 30$ days, the same as the 60-day window's 30.
\textbf{8.} 1.267, a reduction of 0.7\%.
\textbf{9.} The fitted $Q$ is ten times larger (0.0026 against $H = 1.14$ in dollar units) and the filter hedges 1.4\% worse than the window.
\textbf{10.} To about $\pm 0.2$: the filtered ratio's standard deviation is about 0.10.
\textbf{11.} $\times 10$: 0.6\% more error than the window, half-response 9.6 days; $\times 100$: 6.2\% more, 3.0 days.
\textbf{12.} About fifty days for the likelihood's filter (50) and for the window (49); eight days with 100 times the $Q$ (18 with ten times).
\textbf{13.} From $H = 1$, $Q = 0.01$ it reaches $-4\,524$ after ten steps and $-4\,502$ after forty, against the maximum $-4\,470$: the likelihood is flat along a ridge
trading $Q$ against $H$, where EM's steps are short.
\textbf{14.} Filtered 1.20, rolling 1.28 (22 September 2026).
\textbf{15.} When shifts are large and rare compared with the day's noise (a structural break), which a random walk with Gaussian steps does not describe; a model with
occasional jumps in the ratio, or a change-point detector, would react fast only when needed.
\textbf{16.} Only for the uncertainty band, the principled tuning and the handling of missing days; not for the hedge error, which is almost the same.
\textbf{17.} Synchronising the prices (the two spot assessments are taken at different times), adding WTI's lagged change, and hedging with futures rather than
spot prices.
\textbf{18.} Because it uses future data: any backtest with it hedges with hindsight.
\textbf{19.} Named result: \emph{the hedge ratio that moved}: on 2010--2026 data the maximum-likelihood Kalman hedge ratio reduces the hedge-error variance by
0.7\% against the 60-day rolling regression, with an estimated \omterm{def:qm:state-space-models-and-the-kalman-filter:snr}{signal-to-noise ratio} of 0.00047, which makes the filter as slow as the window.
\textbf{20.} The ratio of the state's noise to the observation noise, which the likelihood estimates from the data.
\end{solution}

\begin{solution}{iq:qm:state-space-models-and-the-kalman-filter:1}
You believe something you cannot see, like a hedge ratio, drifts slowly; each day you predict it, see how wrong the prediction of today's prices was, and move your
belief toward the evidence by a fraction that depends on how noisy the prices are compared with how fast the hidden thing moves. You also track how uncertain your
belief is. It is the optimal way to do this when everything is linear and Gaussian.
\iqlookfor{Predict, compare, correct in proportion to relative noise, and carry the uncertainty.}
\end{solution}

\begin{solution}{iq:qm:state-space-models-and-the-kalman-filter:2}
Prior $\alpha \sim \mathcal N(a, P)$, observation $y = z\alpha + \varepsilon$, $\varepsilon \sim \mathcal N(0, H)$. $(\alpha, y)$ are jointly normal with $\Cov = Pz$ and $\Var(y) = z^2P + H$.
Conditioning: $a' = a + \frac{Pz}{z^2P + H}(y - za)$ and $P' = P - \frac{P^2z^2}{z^2P + H}$; the gain is $K = Pz/(z^2P + H)$.
\iqlookfor{The joint-normal argument, not a memorised formula.}
\end{solution}

\begin{solution}{iq:qm:state-space-models-and-the-kalman-filter:3}
Estimate it with the observation noise by maximum likelihood from the \omterm{def:qm:state-space-models-and-the-kalman-filter:ped}{prediction-error decomposition}, then check the hedge out of sample and the speed-noise trade-off;
choosing it by hand to ``react in days'' buys noise. Scale the observation noise with volatility first, or the fit is dominated by volatile periods.
\iqlookfor{Likelihood-based tuning, the trade-off, and heteroskedasticity.}
\end{solution}

\begin{solution}{iq:qm:state-space-models-and-the-kalman-filter:4}
Filtering conditions on data up to $t$; smoothing on all data, past and future. The smoothed path is better for describing history and wrong for trading
simulations, where it is look-ahead bias.
\iqlookfor{Information sets and look-ahead bias.}
\end{solution}

\begin{solution}{iq:qm:state-space-models-and-the-kalman-filter:5}
When the model is nonlinear or non-Gaussian: stochastic volatility, jumps, \omterm{def:qm:robust-statistics-and-heavy-tails:heavy}{heavy-tailed} noise, discrete regimes. It can degenerate (few particles carry the
weight) on surprising observations, its likelihood is noisy (which complicates parameter estimation), and its cost grows with the state dimension.
\iqlookfor{Nonlinearity as the reason, degeneracy and Monte Carlo noise as the risks, \omterm{def:qm:bayesian-methods:ess}{effective sample size} as the monitor.}
\end{solution}

\begin{solution}{iq:qm:state-space-models-and-the-kalman-filter:6}
For any law $q$ of the hidden variables, $\ln p(y \mid \theta) = \E_q[\ln p(y, \alpha \mid \theta)] - \E_q[\ln q] + \mathrm{KL}(q\,\|\,p(\alpha \mid y, \theta))$. The E-step sets $q$ to the
posterior at $\theta_k$, making the KL term zero, so the bound touches the likelihood there; the M-step raises the bound; the likelihood at $\theta_{k+1}$ is at least
the raised bound, hence at least the likelihood at $\theta_k$.
\iqlookfor{The lower bound and the KL gap.}
\end{solution}
